\section{相关技术基础}
本章用于测试正文中的公式、参考文献、表格标题和表格交叉引用。样例内容参考本科论文中的MoE视觉语言模型压缩场景，但只保留最小可检查单元。

\subsection{混合专家模型}
混合专家模型（Mixture of Experts，MoE）通过门控网络为每个输入token选择少量专家，从而在扩大总参数量的同时控制单次推理计算量。早期稀疏门控MoE证明了条件计算在大规模神经网络中的可扩展性\cite{shazeer2017outrageously}，Switch Transformer进一步简化了路由策略\cite{fedus2022switch}。

设第$l$层包含$N$个专家，门控网络给出的权重为$g_i^{(l)}(\bm{x})$，专家输出为$E_i^{(l)}(\bm{x})$，则MoE层输出可写为\eqnref{eq:test_moe_output}。
\begin{equation}
    \label{eq:test_moe_output}
    y^{(l)}(\bm{x}) = \sum_{i=1}^{N} g_i^{(l)}(\bm{x}) E_i^{(l)}(\bm{x})
\end{equation}

\subsection{专家合并与剪枝}
MoE压缩通常分为专家合并和专家剪枝两类。专家合并将多个功能相近的专家融合为一个新专家，专家剪枝则直接删除重要性较低的专家。HC-SMoE、MC-SMoE和REAP分别代表了层次聚类合并、路由感知合并和一次性专家剪枝三类思路\cite{chi2022hcsmoe,li2024mcsmoe,chen2024reap}。

\tabref{chart:merge_comparison}给出三种方法的简化对比。该表用于检查三线表、较长文本列、表题编号和表内英文连字符的Word转换效果。

\begin{table}[!ht]
    \centering
    \caption{三种MoE专家压缩方法对比}
    \label{chart:merge_comparison}
    \fitwidetable{%
    \begin{tabular}{lccc}
    \toprule
        方法 & 聚类或排序依据 & 压缩策略 & 是否需要校准数据 \\ \midrule
        HC-SMoE & 输出响应距离 & 等权平均 & 是 \\
        MergeMoE & 输出L2距离 & 等权平均 & 是 \\
        MC-SMoE & 输出距离+路由频率 & 频率加权平均+低秩补偿 & 是 \\
        REAP & 门控权重+激活范数 & 删除低重要性专家 & 是 \\
    \bottomrule
    \end{tabular}%
    }
\end{table}

\subsection{本章小结}
本章提供了公式、引用和普通数据表测试。后续章节继续放置图片、算法、多列表格和多子图，用于检查DOCX导出后的综合排版效果。
